\section[Retorno de hoja, incidencia areal--volumétrica y orientación constitutiva]{Retorno de hoja, incidencia areal--volumétrica\\y orientación constitutiva}
\label{iii:sec:hojas}
\label{nuclear:constitutiva-antecedente}

La estructura de hojas interviene antes de la evaluación de las magnitudes
electromagnéticas. Su función consiste en distinguir el cierre de una proyección,
la restitución de una orientación y la evolución del registro transportado.
Estas operaciones actúan sobre objetos relacionados, pero de distinto tipo: una
trayectoria proyectada puede haberse cerrado mientras su levantamiento permanece
en la hoja opuesta; la orientación puede haberse restituido mientras la memoria
conserva las dos vueltas efectuadas. La construcción constitutiva debe recibir
esta información con la estructura que la determina.

\subsection{Retorno proyectado y restitución de la cubierta orientada}

Sea $\ell$ un lazo de retorno de una extensión superviviente y sea
$\varepsilon:\langle\ell\rangle\to\{+1,-1\}$ el carácter de su
levantamiento, con $\varepsilon(\ell)=-1$. El fibrado de intervalos asociado
tiene la presentación
\begin{equation}
 \mathcal M=([0,1]\times[-1,1])/{(0,u)\sim(1,-u)}.
 \label{iii:eq:mobius}
\end{equation}
El dato decisivo es el signo del transporte sobre el lazo especificado. La
identificación de los extremos realiza geométricamente ese carácter.

\begin{proposition}[Dos órdenes de clausura]
Una vuelta del lazo cierra la base e invierte la orientación de la fibra. Dos
vueltas restituyen la orientación. En la parametrización angular del lector de
borde, estos retornos corresponden, respectivamente, al cierre areal $2\pi$ y
al cierre volumétrico $4\pi$ de la cubierta orientada.
\end{proposition}
\begin{proof}
La relación de equivalencia de \eqref{iii:eq:mobius} identifica la fibra terminal
con la inicial mediante $u\mapsto-u$. La composición de dos transportes es
$u\mapsto-(-u)=u$. Parametrizar una vuelta de la base por $2\pi$ asigna $4\pi$
a su cuadrado. La afirmación se refiere a este levantamiento y a esta lectura.
\end{proof}

La memoria del TPK conserva, además, la trayectoria recorrida y las actualizaciones
del registro. El cuadrado del carácter de signo es la identidad, pero no impone
que se anulen esas actualizaciones. La representación espinorial, cuando se
utilice, deberá transportar este dato mediante su aplicación propia.

\subsection{Descenso de la incidencia desde el calendario trítico}

En el bloque primitivo $(+1,-1,0)$, el TPK actualiza el modo aritmético mediante
$+1:m\mapsto\Sigma$, $-1:m\mapsto\Pi$ y $0:m\mapsto m$. Su órbita cíclica
es $(\Sigma,\Pi,\Pi)$. El lector de hojas asigna $\Sigma$ a la hoja areal y
$\Pi$ a la volumétrica. En el orden de bases $(e_{\rm ar},e_{\rm vol})$, la
incidencia de aristas queda representada por
\begin{equation}
 F_{\rm av}=\begin{pmatrix}0&1\\1&1\end{pmatrix}.
 \label{iii:eq:fav}
\end{equation}
Los tres pares consecutivos son $\Sigma\to\Pi$, $\Pi\to\Pi$ y
$\Pi\to\Sigma$, cada uno con multiplicidad uno. Los calendarios de longitudes
$9$, $27$ y $108$ contienen $3$, $9$ y $36$ copias del bloque primitivo. Sus
matrices de incidencia son, por tanto, esos múltiplos de $F_{\rm av}$.

Esta construcción conserva la diferencia entre frecuencia cronológica y medida
del espacio de continuaciones. La primera asigna pesos $(1/3,2/3)$ al ciclo
de modos. La segunda se define sobre la envolvente simbólica mínima de memoria
uno que admite exactamente los pares anteriores. Su matriz de adyacencia es
$F_{\rm av}$; los recuentos de trayectorias pertenecen a esa envolvente.

\Needspace{10\baselineskip}
\begin{proposition}[Pesos de retorno de las dos hojas]
Sean $F_0=0$, $F_1=1$ y $F_{n+1}=F_n+F_{n-1}$. Para $n\geq1$,
\begin{equation}
 F_{\rm av}^{n}=
 \begin{pmatrix}F_{n-1}&F_n\\F_n&F_{n+1}\end{pmatrix}.
 \label{iii:eq:fav-powers}
\end{equation}
La medida de Parry de esta envolvente asigna a las hojas pesos proporcionales
a $(1,\varphi^2)$, y la razón de recuentos diagonales converge a $\varphi^{-2}$.
\end{proposition}
\begin{proof}
La fórmula de potencias se prueba por inducción utilizando la recurrencia.
El cuadrado de $F_{\rm av}$ es estrictamente positivo. Su vector positivo de
Perron es $(1,\varphi)$, con autovalor $\varphi$, pues
$\varphi^2=\varphi+1$. Al ser simétrica la matriz, los pesos de Parry son
proporcionales a los cuadrados de las componentes de ese vector. La otra raíz
del polinomio característico es $-\varphi^{-1}$; su módulo es estrictamente
menor que $\varphi$, lo que da el límite de los cocientes diagonales.
\end{proof}

La identificación espectral de esta incidencia es posterior a su construcción
desde el calendario. La coordenada regional $\varphi$ conserva su genealogía
coinductiva anterior: el reconocimiento del autovalor no selecciona el calendario.

\subsection{Marcado de la frontera y razón areal--volumétrica}

Sea $\Pi_{\rm HMT}(X_\infty)$ el valor del lector regional de clausura ya
construido. Con los proyectores diagonales de hoja, se define
\begin{equation}
 D_\pi=\Pi_{\rm HMT}(X_\infty)P_{\rm ar}+P_{\rm vol},\qquad
 \Theta_n=
 \frac{\langle e_{\rm ar},D_\pi F_{\rm av}^{n}e_{\rm ar}\rangle}
 {\langle e_{\rm vol},F_{\rm av}^{n}e_{\rm vol}\rangle}.
 \label{iii:eq:marked-return}
\end{equation}
La normalización volumétrica es unitaria; el cierre regional marca una sola vez
el retorno areal. Aplicando \eqref{iii:eq:fav-powers},
\begin{equation}
 \Theta_n=\Pi_{\rm HMT}(X_\infty)\frac{F_{n-1}}{F_{n+1}}
 \longrightarrow\frac{\pi}{\varphi^2}.
 \label{iii:eq:av-limit}
\end{equation}
La razón de retorno queda determinada por la incidencia, la medida y el marcado.
El cociente $2\pi/4\pi$ describe, en cambio, los dos órdenes de clausura de la
cubierta. Las dos relaciones participan de la lectura de hojas sin sustituirse.

\subsection{Transporte de la orientación a los canales constitutivos}

En una representación de dos canales, sea $\mathcal R$ una involución
autoadjunta y $\mathcal S$ un operador unitario de intercambio que satisface
$\mathcal S\mathcal R\mathcal S^{-1}=-\mathcal R$. Las coordenadas angulares
ya construidas se escriben $x=A_{\rm rad}$ e $y=C^*_{\rm rad}$, con $x>|y|$.
Entonces
\begin{equation}
 B=xI+y\mathcal R,\qquad T=e^{-B},\qquad
 \mathcal ST(y)\mathcal S^{-1}=T(-y).
 \label{iii:eq:channel-conjugation}
\end{equation}
En efecto, la conjugación transforma $B(y)$ en $B(-y)$ y conmuta con su serie
exponencial convergente. Los autovalores $q_\pm=e^{-x\mp y}$ se intercambian.
Para fijar el orden de las respuestas se utiliza la cámara $x>y>0$, donde
$0<q_+<q_-<1$. La aplicación $(x,y)\mapsto(x,-y)$ la transporta a la
cámara opuesta $x>-y>0$, donde $0<q_-<q_+<1$. Ambas pertenecen al dominio
contractivo $x>|y|$. En la frontera $y=0$ los dos canales coinciden.
La conjugación intercambia las cámaras, no conserva el orden estricto de
sus etiquetas; sobre los proyectores actúa mediante
$\mathcal SP_\pm\mathcal S^{-1}=P_\mp$, con
$P_\pm=(I\pm\mathcal R)/2$.
Esta aplicación concreta transmite la orientación al operador constitutivo que
se desarrolla a continuación.
